\chapter{Arquitectura propuesta}

\section{Principios de diseno}

La arquitectura se basa en cuatro principios:

\begin{enumerate}
  \item separar codigo, memoria academica y recursos;
  \item mantener los datos crudos inmutables;
  \item representar decisiones mediante contratos estructurados;
  \item ejecutar transformaciones mediante modulos Python deterministas.
\end{enumerate}

\section{Separacion entre agentes y ejecutores}

Los agentes no escriben ni ejecutan codigo arbitrario sobre los datos. Su papel
es proponer configuraciones, interpretar perfiles y decidir la siguiente fase
del flujo. Los ejecutores leen artefactos desde disco, aplican operaciones
deterministas y devuelven rutas, metricas y logs.

\section{Nodos previstos}

El grafo minimo previsto contiene:

\begin{itemize}
  \item supervisor;
  \item agente limpiador;
  \item ejecutor de limpieza;
  \item agente estructurador;
  \item ejecutor de estructuracion;
  \item agente modelador;
  \item ejecutor de modelado;
  \item evaluador;
  \item redactor;
  \item revision humana para trabajos costosos.
\end{itemize}

\section{Estado global}

El estado global de LangGraph no almacenara datasets completos. Debe contener
rutas, resumenes y artefactos:

\begin{itemize}
  \item identificador de ejecucion;
  \item contexto del proyecto;
  \item rutas de datos crudos, perfiles, datos limpios, tensores e informes;
  \item configuraciones de limpieza, estructuracion y modelado;
  \item metricas y evaluacion;
  \item errores, aprobaciones humanas y artefactos generados.
\end{itemize}

En la implementacion inicial se ha definido un modelo canonico `TFMStateModel`
validado con Pydantic y una vista `TFMState` basada en `TypedDict` para su uso
en LangGraph. Esta separacion permite validar contratos de forma estricta antes
de entregar el estado al grafo y, al mismo tiempo, mantener una representacion
ligera y serializable.

El estado inicial del caso CWRU queda fijado en la fase `dataset\_manifest`,
con el nodo siguiente `manifest\_executor`. El directorio crudo se registra como
\path{codigo/data/raw/cwru_bearing/mat}, el canal principal es `DE\_time` y la
frecuencia objetivo es 12 kHz. Las rutas de artefactos posteriores permanecen
como valores nulos hasta que los ejecutores correspondientes las generen.

\section{Contratos estructurados}

Los contratos iniciales incluyen:

\begin{itemize}
  \item `ProjectContext`, que define dominio, maquina, dataset, objetivo,
        canal principal y frecuencia objetivo;
  \item `DatasetManifestRow` y `DatasetManifest`, que formalizan el inventario
        de ficheros crudos y sus etiquetas;
  \item `DatasetProfileSummary`, que resume el perfil sin almacenar senales;
  \item `CleaningConfig`, `StructuringConfig` y `ModelingConfig`, que fijan las
        configuraciones que podran proponer los agentes;
  \item `SupervisorDecision`, `CleaningDecision`, `StructuringDecision`,
        `ModelingDecision`, `EvaluationDecision` y `ReportDecision`, que
        estructuran las respuestas de los agentes;
  \item `ExecutorResult` y resultados especializados, que normalizan las
        salidas de los ejecutores deterministas;
  \item `MetricsReport` y `EvaluationResult`, que estructuran la evaluacion;
  \item `PipelineError`, `HumanApproval` y `ArtifactRef`, que mantienen
        trazabilidad de errores, revisiones y artefactos.
\end{itemize}

Estos esquemas aplican validacion estricta y rechazan campos no declarados. Por
ejemplo, una fila normal del manifiesto debe declarar los campos de fallo como
nulos, mientras que una fila con fallo debe incluir tipo de fallo y diametro. De
forma analoga, un resultado de ejecutor con estado fallido debe incluir un error
estructurado, y una decision no terminal del supervisor debe indicar el siguiente
nodo del grafo.

\section{Trazabilidad}

Cada ejecucion debe permitir responder:

\begin{itemize}
  \item que datos se usaron;
  \item que configuracion fue validada;
  \item que ejecutor produjo cada artefacto;
  \item que metricas se obtuvieron;
  \item que limitaciones permanecen.
\end{itemize}
